\clearpage
\SourceLeaf{611}
\ArithmeticHead{582}{L}{INHALT.}
\begingroup
\newcommand{\ContentsRow}[5]{\par\noindent\hangindent=#1\hspace*{#1}#2\hspace{.4em}\mbox{\footnotesize #3}\hfill\mbox{#4\quad #5}\par}
\noindent\textls[120]{Zur Theorie der unendlichen Reihe} $F(\alpha,\beta,\gamma,x)$.\par
\ContentsRow{1em}{I.\quad Entwicklungen in Kettenbrüche}{.\,.\,.\,.\,.\,.\,.\,.\,.\,.\,.\,.\,.\,.\,.\,.\,.\,.\,.\,.}{Seite}{326}
\ContentsRow{3em}{Bemerkungen}{.\,.\,.\,.\,.\,.\,.\,.\,.\,.\,.\,.\,.\,.\,.\,.\,.\,.\,.\,.\,.\,.\,.\,.}{—}{330}
\ContentsRow{1em}{II.\quad Allgemeines über die Reihe $F(\alpha,\beta,\gamma,x)$}{.\,.\,.\,.\,.\,.\,.\,.\,.\,.\,.\,.\,.\,.\,.\,.\,.}{—}{332}
\ContentsRow{3em}{Bemerkungen}{.\,.\,.\,.\,.\,.\,.\,.\,.\,.\,.\,.\,.\,.\,.\,.\,.\,.\,.\,.\,.\,.\,.\,.}{—}{337}
\ContentsRow{1em}{III.\quad Einiges über die unendliche Reihe $1+\dfrac{\alpha\mathord{.}\beta}{1\mathord{.}\gamma}\,x+\mathord{\cdot}\,\mathord{\cdot}\,\mathord{\cdot}$}{.\,.\,.\,.\,.\,.\,.\,.\,.\,.\,.\,.\,.}{—}{338}
\ContentsRow{3em}{Bemerkungen}{.\,.\,.\,.\,.\,.\,.\,.\,.\,.\,.\,.\,.\,.\,.\,.\,.\,.\,.\,.\,.\,.\,.\,.}{—}{353}
\ContentsRow{1em}{IV.\quad Nachträge zum art.\ 6.\ der Abhandlung von 1812}{.\,.\,.\,.\,.\,.\,.\,.\,.\,.\,.\,.\,.\,.}{—}{355}
\ContentsRow{3em}{Bemerkung}{.\,.\,.\,.\,.\,.\,.\,.\,.\,.\,.\,.\,.\,.\,.\,.\,.\,.\,.\,.\,.\,.\,.\,.\,.}{—}{359}
\par\medskip\noindent\hspace*{1em}\emph{Briefwechsel:}\par
\ContentsRow{2em}{1.\ \textsc{Gauss} an \textsc{Bessel}, 1810 Oktober 21}{.\,.\,.\,.\,.\,.\,.\,.\,.\,.\,.\,.\,.\,.\,.\,.\,.}{—}{360}
\ContentsRow{2em}{2.\ \textsc{Gauss} an \textsc{Olbers}, 1811 Oktober 17}{.\,.\,.\,.\,.\,.\,.\,.\,.\,.\,.\,.\,.\,.\,.\,.\,.}{—}{361}
\ContentsRow{2em}{3.\ \textsc{Gauss} an \textsc{Bessel}, 1811 November 21}{.\,.\,.\,.\,.\,.\,.\,.\,.\,.\,.\,.\,.\,.\,.\,.}{—}{362}
\ContentsRow{2em}{4.\ \textsc{Gauss} an \textsc{Bessel}, 1811 Dezember 18}{.\,.\,.\,.\,.\,.\,.\,.\,.\,.\,.\,.\,.\,.\,.\,.}{—}{365}
\ContentsRow{2em}{5.\ \textsc{Gauss} an \textsc{Laplace}, 1812 Januar 30}{.\,.\,.\,.\,.\,.\,.\,.\,.\,.\,.\,.\,.\,.\,.\,.\,.}{—}{371}
\ContentsRow{2em}{6.\ \textsc{Gauss} an \textsc{Bessel}, 1812 Mai 5}{.\,.\,.\,.\,.\,.\,.\,.\,.\,.\,.\,.\,.\,.\,.\,.\,.\,.}{—}{374}
\ContentsRow{2em}{7.\ \textsc{Gauss} an \textsc{Laplace}, 1812 November 5}{.\,.\,.\,.\,.\,.\,.\,.\,.\,.\,.\,.\,.\,.\,.\,.}{—}{378}
\ContentsRow{2em}{Bemerkungen}{.\,.\,.\,.\,.\,.\,.\,.\,.\,.\,.\,.\,.\,.\,.\,.\,.\,.\,.\,.\,.\,.\,.\,.\,.\,.}{—}{379}
\par\medskip\noindent\textls[120]{Zur Lehre von den Reihen.}\par
\noindent\hspace*{1em}I.\quad Neue Methode die Summe der divergierenden Reihe\par
\[1-1+2-6+24-\text{etc.}=0{,}5963\,\mathord{.}\,\mathord{.}\,\mathord{.}\]
\ContentsRow{2em}{zu finden}{.\,.\,.\,.\,.\,.\,.\,.\,.\,.\,.\,.\,.\,.\,.\,.\,.\,.\,.\,.\,.\,.\,.\,.\,.\,.}{—}{382}
\ContentsRow{3em}{Bemerkungen}{.\,.\,.\,.\,.\,.\,.\,.\,.\,.\,.\,.\,.\,.\,.\,.\,.\,.\,.\,.\,.\,.\,.\,.}{—}{385}
\ContentsRow{1em}{II.\quad Grundbegriffe der Lehre von den Reihen}{.\,.\,.\,.\,.\,.\,.\,.\,.\,.\,.\,.\,.\,.\,.\,.\,.}{—}{390}
\ContentsRow{3em}{Bemerkungen}{.\,.\,.\,.\,.\,.\,.\,.\,.\,.\,.\,.\,.\,.\,.\,.\,.\,.\,.\,.\,.\,.\,.\,.}{—}{395}
\ContentsRow{1em}{III.\quad Fragen zur Metaphysik der Mathematik}{.\,.\,.\,.\,.\,.\,.\,.\,.\,.\,.\,.\,.\,.\,.\,.\,.}{—}{396}
\ContentsRow{3em}{Bemerkung}{.\,.\,.\,.\,.\,.\,.\,.\,.\,.\,.\,.\,.\,.\,.\,.\,.\,.\,.\,.\,.\,.\,.\,.\,.}{—}{397}
\ContentsRow{1em}{IV.\quad Darstellung von diskontinuierlichen Funktionen}{.\,.\,.\,.\,.\,.\,.\,.\,.\,.\,.\,.\,.\,.}{—}{398}
\ContentsRow{3em}{Bemerkung}{.\,.\,.\,.\,.\,.\,.\,.\,.\,.\,.\,.\,.\,.\,.\,.\,.\,.\,.\,.\,.\,.\,.\,.\,.}{—}{399}
\ContentsRow{1em}{V.\quad Convergenz der Reihen, in welche die periodischen Functionen einer veränderlichen Größe entwickelt werden}{.\,.\,.\,.\,.\,.\,.\,.\,.\,.\,.\,.\,.\,.\,.\,.\,.\,.\,.\,.\,.}{—}{400}
\ContentsRow{1em}{VI.\quad Bestimmung der Convergenz der Reihen, in welche die periodischen Functionen einer veränderlichen Größe entwickelt werden}{.\,.\,.\,.\,.\,.\,.\,.\,.\,.\,.\,.\,.\,.\,.}{—}{407}
\ContentsRow{1em}{VII.\quad Über die Konvergenz der Entwicklung der Mittelpunktsgleichung}{.\,.\,.\,.\,.\,.\,.\,.}{—}{420}
\par\medskip\noindent\hspace*{1em}\emph{Briefwechsel:}\par
\ContentsRow{2em}{1.\ \textsc{Gauss} an \textsc{F. Hindenburg}, 1799 Oktober 8}{.\,.\,.\,.\,.\,.\,.\,.\,.\,.\,.\,.\,.\,.}{—}{429}
\ContentsRow{2em}{2.\ \textsc{Gauss} an \textsc{Schumacher}, 1849 Dezember 4}{.\,.\,.\,.\,.\,.\,.\,.\,.\,.\,.\,.\,.\,.\,.}{—}{431}
\ContentsRow{2em}{3.\ \textsc{Gauss} an \textsc{Schumacher}, 1849 Dezember 6}{.\,.\,.\,.\,.\,.\,.\,.\,.\,.\,.\,.\,.\,.\,.}{—}{431}
\ContentsRow{2em}{4.\ \textsc{Gauss} an \textsc{Schumacher}, 1850 Februar 5}{.\,.\,.\,.\,.\,.\,.\,.\,.\,.\,.\,.\,.\,.\,.}{—}{433}
\ContentsRow{2em}{5.\ \textsc{Gauss} an \textsc{Schumacher}, 1850 September 1}{.\,.\,.\,.\,.\,.\,.\,.\,.\,.\,.\,.\,.\,.}{—}{434}
\ContentsRow{2em}{6.\ \textsc{Gauss} an \textsc{H. G. Grassmann}, 1844 Dezember 14}{.\,.\,.\,.\,.\,.\,.\,.\,.\,.\,.\,.}{—}{436}
\par\smallskip\noindent\hspace*{1em}Bemerkungen zu den Abschnitten V, VI, VII und zum Briefwechsel.\par
\ContentsRow{2em}{Die Abschnitte V und VI}{.\,.\,.\,.\,.\,.\,.\,.\,.\,.\,.\,.\,.\,.\,.\,.\,.\,.\,.\,.\,.\,.}{—}{437}
\endgroup
